\section{Contingent Collaborative Tasks}
\label{sec:contingent_collaborative_tasks}

This section develops the high-level coordination layer for contingent
collaborative tasks. The main objective is to determine how each robot
updates its local task and adapts its discrete plan when contingent
requests are exchanged during runtime. Two request types are considered,
namely service requests and formation requests. The focus is on logical
specification, event monitoring, communication, and discrete plan
adaptation. The continuous realization of navigation and formation
maintenance is only treated as an execution interface and is not
developed here.

\subsection{RTL and sc-RTL preliminaries}
\label{subsec:rtl_scrtl_prelim}

The contingent-task formulation in this section relies on regular
temporal logic and, in particular, on its syntactically co-safe
fragment. A brief introduction is therefore required before the local
task in \eqref{eq:task_contingent} is defined.

Let $AP$ be a finite set of atomic propositions. A word over $AP$ is an
element of $(2^{AP})^\omega$ if it is infinite, and an element of
$(2^{AP})^\ast$ if it is finite. In the present setting, each symbol in
such a word records which propositions are true at a given discrete
observation instant.

\begin{definition}[RTL specification]
\label{def:rtl_specification}
A regular temporal logic formula over $AP$ is a temporal specification
that defines a language over words in $(2^{AP})^\omega$ or
$(2^{AP})^\ast$. It is used to describe sequential and temporal
requirements over proposition traces generated by robot motions.

For a robot motion trace, satisfaction of an RTL formula means that the
induced proposition word belongs to the language specified by that
formula.
\end{definition}

In this chapter, the static local task is allowed to be recurrent, hence
it is modeled as an RTL formula over infinite executions. By contrast,
the contingent requests introduced later are short-horizon tasks that
must be completed in finite time. This motivates the use of the
syntactically co-safe fragment.

\begin{definition}[sc-RTL formula]
\label{def:scrtl_formula}
A syntactically co-safe RTL formula, abbreviated as sc-RTL, is an RTL
formula whose satisfaction can be certified by a finite prefix of the
generated proposition word.

Equivalently, for every sc-RTL formula $\psi$, if a trace satisfies
$\psi$, then there exists a finite word prefix after which satisfaction
of $\psi$ is already determined.
\end{definition}

Definition~\ref{def:scrtl_formula} is precisely the property needed for
contingent collaborative tasks. Indeed, a service request and a
formation request should both be completed after a finite execution
segment, rather than requiring an infinite recurrent behavior.

For each sc-RTL formula $\psi$, let $[\psi]$ denote its equivalent
co-safe temporal specification over proposition words. This notation is
used only to interface task formulas with finite-state automata. In
particular, $[\psi]$ admits a finite accepting automaton whose accepting
states characterize completion of the corresponding contingent task.
This observation will be used in
Subsection~\ref{subsec:event_based_adaptation} for event-based plan
adaptation.

The static local specification of robot $i$ is denoted by
$\varphi_i^{\texttt{l}}$. The contingent service request
$\varphi_{ji,t}^{\texttt{s}}$ and the contingent formation request
$\varphi_{ij,t}^{\texttt{f}}$ are both taken to be sc-RTL formulas.
Hence they can be recognized by finite automata and can be inserted into
the online planning process as finite tasks with explicit completion
events.

\subsection{Problem formulation with contingent collaborative tasks}
\label{subsec:problem_contingent_tasks}

Consider the robot set
$\mathcal{N}\triangleq \{1,\cdots,N\}$.
For each robot $i\in\mathcal{N}$, let $\mathcal{N}_i$ denote the set of
its communication neighbors. The workspace and the region propositions
are inherited from the previous sections. In particular,
$R\triangleq \{R_1,\cdots,R_M\}$ is the set of atomic propositions
associated with the regions in $\Pi$.

\begin{definition}[Contingent collaborative requests]
\label{def:contingent_requests}
A \emph{service request} from robot $j$ to robot $i$ at time $t$ is a
feasible sc-RTL formula $\varphi_{ji,t}^{\texttt{s}}$ over $R$.

A \emph{formation request} from robot $i$ to robot $j$ at time $t$
consists of a desired relative deployment descriptor together with a
feasible sc-RTL formula $\varphi_{ij,t}^{\texttt{f}}$ over $R$.
\end{definition}

The local task of each robot keeps the original three-part structure.
To express the contingent parts, the following auxiliary propositions are
introduced.

\begin{definition}[Auxiliary propositions]
\label{def:auxiliary_props}
For each robot $i\in\mathcal{N}$, define
\[
O_i^{\texttt{s}}\triangleq \{o_{ji}^{\texttt{s}} \,:\, j\in\mathcal{N}_i\},
\qquad
O_i^{\texttt{f}}\triangleq \{o_{ij}^{\texttt{f}} \,:\, j\in\mathcal{N}_i\},
\]
where $o_{ji}^{\texttt{s}}(t)=\top$ when robot $i$ confirms the service
request received from robot $j$ at time $t$, and
$o_{ij}^{\texttt{f}}(t)=\top$ when robot $j$ confirms the formation
request sent by robot $i$ at time $t$.

Define also
\[
H_i\triangleq \{h_{ij}\,:\, j\in\mathcal{N}_i\},
\qquad
Z_i\triangleq \{z_{ij}\,:\, j\in\mathcal{N}_i\},
\]
where $h_{ij}(t)=\top$ means that the requested formation associated with
robot pair $i,j$ is currently maintained by the execution layer, and
$z_{ij}(t)=\top$ means that robot $i$ releases robot $j$ from that
formation request.
\end{definition}

The complete proposition set of robot $i$ is therefore
\begin{equation}
\label{eq:ap_i_contingent}
AP_i \triangleq R \cup O_i^{\texttt{s}} \cup O_i^{\texttt{f}}
\cup H_i \cup Z_i,
\end{equation}
where the symbols are given in Definition~\ref{def:auxiliary_props}.

Each robot $i$ is assigned a local RTL task over $AP_i$ with the
following structure:
\begin{equation}
\label{eq:task_contingent}
\varphi_i \triangleq
\varphi_i^{\texttt{l}}
\wedge
\varphi_i^{\texttt{s}}
\wedge
\varphi_i^{\texttt{f}},
\end{equation}
where $\varphi_i^{\texttt{l}}$ is the static local task over $R$,
$\varphi_i^{\texttt{s}}$ captures contingent service requests, and
$\varphi_i^{\texttt{f}}$ captures contingent formation requests.

The service part is defined by
\begin{equation}
\label{eq:service_task_contingent}
\varphi_i^{\texttt{s}} \triangleq
\bigwedge_{j\in\mathcal{N}_i}
\square
\big(
o_{ji}^{\texttt{s}} \rightarrow \varphi_{ji,t}^{\texttt{s}}
\big),
\end{equation}
which means that once robot $i$ confirms a service request sent by robot
$j$, the corresponding contingent task must be fulfilled afterwards.

The formation part is defined by
\begin{equation}
\label{eq:formation_task_contingent}
\varphi_i^{\texttt{f}} \triangleq
\bigwedge_{j\in\mathcal{N}_i}
\square
\Big(
o_{ij}^{\texttt{f}} \rightarrow
\big(
\varphi_{ij,t}^{\texttt{f}} \wedge (h_{ij}\,\mathsf{U}\,z_{ij})
\big)
\Big),
\end{equation}
which means that once robot $j$ confirms a formation request sent by
robot $i$, robot $i$ must complete the associated short-term task
$\varphi_{ij,t}^{\texttt{f}}$, while the requested formation is
maintained until the release signal $z_{ij}$ is issued.

Equation~\eqref{eq:task_contingent} is the operative task formula of
robot $i$. It should be emphasized that, unlike the standard local-task
setting in multi-robot temporal-logic planning
\cite{bhatia2010sampling,guo2015multi}, the full content of
$\varphi_i$ cannot be known offline. This is because the truth values of
$o_{ji}^{\texttt{s}}$, $o_{ij}^{\texttt{f}}$, $h_{ij}$, and $z_{ij}$
are determined only during runtime.

\begin{problem}
\label{prob:contingent_collaborative}
Consider a team of robots together with their motion abstractions and
the local RTL task in \eqref{eq:task_contingent}. The objective is to
synthesize, for each robot $i\in\mathcal{N}$, a distributed event-driven
coordination and planning scheme such that $\varphi_i$ is satisfied,
while contingent service and formation requests may be exchanged during
runtime.
\end{problem}

\subsection{Triggering events and communication protocol}
\label{subsec:trigger_protocol}

The contingent structure in \eqref{eq:task_contingent} requires a
runtime mechanism for event monitoring and request handling. This part
follows the same logic as the original construction and separates
triggering events from the request-exchange protocol.

\subsubsection*{Triggering events}

\begin{definition}[Triggering events]
\label{def:triggering_events}
For each robot $i\in\mathcal{N}$, the following events are monitored.

\emph{Region-crossing event}. This event occurs when robot $i$ enters or
leaves a region $\pi_\ell\in\Pi$. It determines the truth value of the
proposition $R_\ell\in R$.

\emph{Request-reply event}. This event occurs when robot $i$ sends a
service or formation request to a neighbor and receives an immediate
reply. If the request is confirmed, the associated proposition
$o_{ji}^{\texttt{s}}$ or $o_{ij}^{\texttt{f}}$ becomes true for a short
time interval of length $\delta_d>0$.

\emph{Service-finish event}. This event occurs when a confirmed service
request $\varphi_{ji,t}^{\texttt{s}}$ has been satisfied by the trace of
robot $i$ after time $t$.

\emph{Formation-finish event}. This event occurs when the short-term
formation task $\varphi_{ij,t}^{\texttt{f}}$ associated with a confirmed
formation request has been satisfied by robot $i$. At that instant the
release signal $z_{ij}$ is issued.
\end{definition}

The first two events can be monitored directly from region membership and
message exchange. The last two events depend on the current state of the
task automaton. Their detection will therefore be tied to the plan
adaptation mechanism in Subsection~\ref{subsec:event_based_adaptation}.

\subsubsection*{Communication protocol}

The protocol keeps track of which contingent requests are currently being
served. For each pair $j\in\mathcal{N}_i$, let
$r_{ji}^{\texttt{s}}(t)\in\mathbb{B}$ indicate whether robot $i$ is
serving a service request from robot $j$ at time $t$, and let
$r_{ji}^{\texttt{f}}(t)\in\mathbb{B}$ indicate whether robot $i$ is
serving a formation request associated with robot $j$ at time $t$.
In addition, let $r_{ii}^{\texttt{f}}(t)\in\mathbb{B}$ indicate whether
robot $i$ is currently executing the short-term task induced by one of
its own confirmed formation requests.

Initially,
$r_{ji}^{\texttt{s}}(0)=\bot$,
$r_{ji}^{\texttt{f}}(0)=\bot$,
and
$r_{ii}^{\texttt{f}}(0)=\bot$.

A service request $\varphi_{ij,t}^{\texttt{s}}$ can be sent by robot $i$
to robot $j$ only when robot $j$ is not still serving the previous
service request from robot $i$, namely when
$r_{ij}^{\texttt{s}}(t)=\bot$.
A formation request can be sent by robot $i$ to robot $j$ only when
robot $j$ is not still serving the previous formation request from
robot $i$, namely when $r_{ij}^{\texttt{f}}(t)=\bot$, and robot $i$ is
not currently executing an active formation task of its own, namely when
$r_{ii}^{\texttt{f}}(t)=\bot$.

When robot $i$ receives a request from robot $j$ at time $t$, it first
checks whether it is currently involved in an active formation mode. If
$r_{ii}^{\texttt{f}}(t)=\top$ or $r_{ji}^{\texttt{f}}(t)=\top$ for some
$j\in\mathcal{N}_i$, the request is denied. Otherwise the request is
confirmed and the corresponding proposition
$o_{ji}^{\texttt{s}}$ or $o_{ij}^{\texttt{f}}$ is updated according to
Definition~\ref{def:triggering_events}.

The above rules imply the following structural property.

\begin{lemma}
\label{lem:request_capacity}
Each robot may confirm multiple service requests in sequence, but at
most one formation request can remain active at a time.
\end{lemma}

\begin{proof}
The protocol allows a service request whenever the previous service
request from the same requester has already been discharged. By
contrast, any active formation mode blocks further request acceptance.
Hence multiple service requests can be accumulated through time, whereas
formation requests are mutually exclusive.
\end{proof}

\subsection{Discrete plan synthesis}
\label{subsec:discrete_synthesis}

At time $t=0$, no contingent request has yet been confirmed. Hence only
the static task $\varphi_i^{\texttt{l}}$ is active initially. The first
step is therefore the same as in the static setting of the previous
chapter.

\textbf{Motion abstraction.}
The motion of robot $i$ is abstracted by a weighted finite transition
system
\[
\mathcal{T}_i \triangleq
(\Pi_i,\longrightarrow,R,L,\Pi_{i,0},W),
\]
where $\Pi_i\triangleq \{\pi_{i,0}\}\cup\Pi$,
$\pi_{i,0}\triangleq \mathcal{B}(p_i(0),0)$ is the initial state,
$L(\pi_\ell)=R_\ell$, and the transition weight
$W(\pi_m,\pi_n)$ is the travel cost between the corresponding regions.

\textbf{Task automaton.}
Let $\mathcal{B}_{i,s}$ be the NBA associated with the LTL abstraction
$[\varphi_i^{\texttt{l}}]$ of the static RTL formula
$\varphi_i^{\texttt{l}}$ \cite{reynolds2001continuous,
gastin2001fast,baier2008principles}.

\textbf{Product construction.}
The initial weighted product automaton is
\begin{equation}
\label{eq:initial_product_contingent}
\mathcal{A}_{p,i}\triangleq
\mathcal{T}_i \times \mathcal{B}_{i,s},
\end{equation}
where the notation is the same as in the previous chapter.

\textbf{Plan extraction.}
An accepting run of $\mathcal{A}_{p,i}$ with minimum prefix-suffix cost
is computed using the standard graph-search procedure
\cite{guo2012infeasible}. Its projection onto $\Pi_i$ yields the initial
discrete plan
\begin{equation}
\label{eq:initial_plan_contingent}
\tau_{i,0} \triangleq
\pi_{i0}\pi_{i1}\cdots
\big(
\pi_{ik_i}\pi_{i(k_i+1)}\cdots\pi_{iK_i}
\big)^\omega,
\end{equation}
where the prefix is finite and the suffix is repeated infinitely often.

The correctness of the initial plan follows from the same argument as in
the static case.

\begin{theorem}
\label{thm:initial_plan_contingent}
If no contingent request is confirmed by robot $i$, then the trajectory
generated by executing $\tau_{i,0}$ satisfies the local task
$\varphi_i$.
\end{theorem}

\begin{proof}
If no contingent request is confirmed, then
$o_{ji}^{\texttt{s}}(t)=\bot$ and $o_{ij}^{\texttt{f}}(t)=\bot$ for all
$j\in\mathcal{N}_i$ and all $t\geq 0$. Hence the contingent parts
\eqref{eq:service_task_contingent} and
\eqref{eq:formation_task_contingent} are vacuously satisfied, and
$\varphi_i$ reduces to $\varphi_i^{\texttt{l}}$. The claim then follows
from the standard correctness of the product-automaton construction.
\end{proof}

\begin{algorithm}[t]
\caption{Event-driven discrete synthesis and adaptation for robot $i$}
\label{alg:contingent_planning}
\KwIn{$\mathcal{T}_i$, $\varphi_i^{\texttt{l}}$, event stream}
\KwOut{Discrete plan $\tau_i$}

Construct $\mathcal{B}_{i,s}$ for $[\varphi_i^{\texttt{l}}]$\;
Construct $\mathcal{A}_{p,i}$ by \eqref{eq:initial_product_contingent}\;
Compute an accepting run and project it to $\tau_{i,0}$\;
Set $\tau_i\leftarrow \tau_{i,0}$\;

\While{robot $i$ is active}{
  Execute the current plan until the next triggering event\;

  \If{a request-reply event confirms a new contingent task}{
    update the intersection automaton $\mathcal{A}_{[\varphi_i]}$\;
    reconstruct the product automaton from the current product state\;
    recompute an accepting run and update $\tau_i$\;
  }

  \If{a service-finish event occurs}{
    remove the completed service request from the active request set\;
    update $\mathcal{A}_{[\varphi_i]}$ and recompute $\tau_i$\;
  }

  \If{a formation-finish event occurs}{
    send the release signal $z_{ij}$\;
    remove the completed formation task from the active request set\;
    update $\mathcal{A}_{[\varphi_i]}$ and recompute $\tau_i$\;
  }

  \If{a region-crossing event occurs}{
    update the current discrete state and the next target region\;
  }
}
\end{algorithm}

\subsection{Event-based plan adaptation}
\label{subsec:event_based_adaptation}

The initial plan from Subsection~\ref{subsec:discrete_synthesis} is
valid only before contingent requests are confirmed. Once a new request
is accepted, robot $i$ must update the task automaton while preserving
the progress already made on the static task and on the previously
confirmed contingent tasks.

The adaptation procedure is organized into four cases.

\textbf{Case I.}
Robot $i$ confirms its first service request
$\varphi_{ji,t_j}^{\texttt{s}}$ from neighbor $j$ at time $t_j>0$.

Let $\mathcal{B}_{[\varphi_{ji,t_j}^{\texttt{s}}]}$ be the NBA
associated with $[\varphi_{ji,t_j}^{\texttt{s}}]$. Let
$q_{s,t_j}$ be the current state of the static-task automaton obtained
from the current product state at time $t_j$.

\begin{definition}[Single-request prioritized intersection]
\label{def:single_intersection}
The intersection of
$\mathcal{B}_{[\varphi_{ji,t_j}^{\texttt{s}}]}$ and
$\mathcal{B}_{[\varphi_i^{\texttt{l}}]}$ is the NBA
\begin{equation}
\label{eq:single_intersection}
\mathcal{A}_{[\varphi_i]} \triangleq
(Q,2^R,\delta,Q_0,F),
\end{equation}
where
$Q\triangleq Q_j\times Q_s\times\{1,2\}$,
$Q_0\triangleq Q_{j,0}\times\{q_{s,t_j}\}\times\{1\}$,
and
$F\triangleq F_j\times F_s\times\{2\}$.

A transition
$\langle \check q_j,\check q_s,\check c\rangle
\in
\delta(\langle q_j,q_s,c\rangle,l)$
is enabled if
$\check q_j\in\delta_j(q_j,l)$,
$\check q_s\in\delta_s(q_s,l)$,
and one of the following conditions holds:
$c=1$, $q_j\notin F_j$, and $\check c=1$;
or
$c=1$, $q_j\in F_j$, and $\check c=2$;
or
$c=2$ and $\check c=2$.
\end{definition}

Definition~\ref{def:single_intersection} gives the service request
priority over the repeating suffix of the static task. The automaton
remains in the first layer until the request automaton reaches its
accepting set, and only then continues in the second layer to recover
the recurrent static behavior.

\begin{lemma}
\label{lem:single_intersection_correct}
If there exists an $\omega$-word $w\in (2^R)^\omega$ such that
$w\models[\varphi_{ji,t_j}^{\texttt{s}}]$ and
$w\models[\varphi_i^{\texttt{l}}]$, then the automaton
$\mathcal{A}_{[\varphi_i]}$ in \eqref{eq:single_intersection} has at
least one accepting run.
\end{lemma}

\begin{proof}
Any word satisfying the co-safe service request reaches an accepting
state of its request automaton in finite time. The same word also yields
an accepting run of the static-task automaton. The two-layer
construction then concatenates the finite accepting prefix of the
request automaton with the accepting suffix of the static-task
automaton, which produces an accepting run of
$\mathcal{A}_{[\varphi_i]}$.
\end{proof}

\textbf{Case II.}
Robot $i$ confirms a new service request while several earlier service
requests are still active.

Assume that the active service requests are
$\varphi_{1i,t_1}^{\texttt{s}},\cdots,
\varphi_{mi,t_m}^{\texttt{s}}$,
ordered according to their confirmation times. Let
$\mathcal{B}_{[\varphi_{gi,t_g}^{\texttt{s}}]}$ be the NBA associated
with each request, and let $q_{g,t_m}$ and $q_{s,t_m}$ denote the
current automaton states extracted from the current product state at
time $t_m$.

\begin{definition}[Layered prioritized intersection]
\label{def:layered_intersection}
The intersection of
$\mathcal{B}_{[\varphi_{gi,t_g}^{\texttt{s}}]}$,
$g=1,\cdots,m$,
and $\mathcal{B}_{[\varphi_i^{\texttt{l}}]}$ is the NBA
\begin{equation}
\label{eq:layered_intersection}
\mathcal{A}_{[\varphi_i]} \triangleq
(Q,2^R,\delta,Q_0,\mathcal F),
\end{equation}
where
\begin{equation}
\label{eq:layered_intersection_state}
Q \triangleq
Q_1\times\cdots\times Q_m\times Q_s\times\{1,\cdots,m+1\},
\end{equation}
\begin{equation}
\label{eq:layered_intersection_init}
Q_0 \triangleq
\{q_{1,t_m}\}\times\cdots\times\{q_{m-1,t_m}\}
\times Q_{m,0}\times\{q_{s,t_m}\}\times\{1\},
\end{equation}
and
\begin{equation}
\label{eq:layered_intersection_accept}
\mathcal F \triangleq
F_1\times\cdots\times F_m\times F_s\times\{m+1\}.
\end{equation}

A transition
$\langle \check q_1,\cdots,\check q_m,\check q_s,\check c\rangle$
is enabled if
$\check q_g\in\delta_g(q_g,l)$ for all $g=1,\cdots,m$,
$\check q_s\in\delta_s(q_s,l)$, and
the layer update satisfies one of the following rules:
$q_c\notin F_c$ and $\check c=c$;
or
$q_c\in F_c$ and $\check c=c+1$;
or
$c=m+1$ and $\check c=m+1$.
\end{definition}

The layered construction in
Definition~\ref{def:layered_intersection} preserves the current progress
of all previously confirmed requests. It also records the order in which
the contingent tasks must be discharged before the static suffix resumes.

\begin{lemma}
\label{lem:layered_intersection_correct}
If there exists an $\omega$-word $w\in (2^R)^\omega$ that satisfies all
active service requests and also satisfies $[\varphi_i^{\texttt{l}}]$,
then the automaton in \eqref{eq:layered_intersection} has at least one
accepting run.
\end{lemma}

\begin{proof}
The proof follows the same argument as
Lemma~\ref{lem:single_intersection_correct}. Each co-safe request is
completed in finite time. The layered index records these completions in
sequence. Once the last request layer is passed, the run remains on the
accepting suffix of the static-task automaton.
\end{proof}

\textbf{Case III.}
Robot $i$ receives a new service request from a neighbor whose previous
service request has already been completed.

In this case, the completed request is removed from the active set, the
new request is inserted as the latest layer, and the intersection
automaton is reconstructed by
Definition~\ref{def:layered_intersection} from the current product
state. This update preserves the progress on all other active requests
and on the static task.

\textbf{Case IV.}
A formation request sent by robot $i$ is confirmed by neighbor $j$.

The associated formation task $\varphi_{ij,t_j}^{\texttt{f}}$ is an
sc-RTL formula and is treated in the same manner as a newly confirmed
service request at the discrete-planning level. More precisely, the NBA
associated with $[\varphi_{ij,t_j}^{\texttt{f}}]$ is added as a new
layer in the intersection automaton by the same construction as in
Definition~\ref{def:layered_intersection}. Once the accepting set of
that layer is reached, a formation-finish event is detected and the
release signal $z_{ij}$ is generated.

The foregoing cases lead directly to request completion detection.

\begin{theorem}
\label{thm:request_completion}
Assume that every confirmed contingent request is feasible. Then each
confirmed service request $\varphi_{ji,t}^{\texttt{s}}$ and each
confirmed formation task $\varphi_{ij,t}^{\texttt{f}}$ is completed in
finite time, provided that the updated product automaton admits an
accepting run after each event-triggered reconstruction.
\end{theorem}

\begin{proof}
By construction, every active contingent request occupies a dedicated
layer in the current intersection automaton. Since the request formulas
are co-safe, reaching the corresponding accepting set requires only a
finite prefix of the trace. If the updated product automaton admits an
accepting run, that accepting run must eventually traverse each active
request layer and hence complete each contingent task in finite time.
\end{proof}

\subsection{Overall structure}
\label{subsec:overall_structure_contingent}

The overall coordination architecture follows the same modular logic as
the original development, but the present section only retains the
high-level layers.

\textbf{Communication module.}
This module exchanges contingent service requests and formation requests
with the neighbors. It also receives request confirmations and release
signals. Once a request is confirmed, the corresponding contingent task
is forwarded to the discrete planning module.

\textbf{Event monitoring module.}
This module detects region-crossing events and request-reply events
directly. It also detects service-finish events and formation-finish
events by observing the current automaton state in the discrete planning
module.

\textbf{Discrete planning module.}
This module computes the initial plan from
\eqref{eq:initial_product_contingent} and updates it whenever the task
automaton changes. In particular, it reconstructs the prioritized
intersection automaton, rebuilds the product automaton from the current
product state, and computes a new accepting run according to
Algorithm~\ref{alg:contingent_planning}.

\textbf{Execution module.}
This module receives the next goal region from the discrete planner and
executes the corresponding motion segment. During an active formation
request, it also maintains the requested formation until the release
signal is issued. The detailed controller design is deferred to the
later execution-oriented part of the chapter.

The modular decomposition above yields the main correctness statement.

\begin{theorem}
\label{thm:overall_correctness_contingent}
Assume that all confirmed contingent service and formation tasks are
feasible, and assume that each event-triggered product automaton admits
an accepting run. Then the modular coordination scheme described above
generates, for each robot $i\in\mathcal{N}$, an execution that satisfies
the local task formula $\varphi_i$ in \eqref{eq:task_contingent}.
\end{theorem}

\begin{proof}
The static part $\varphi_i^{\texttt{l}}$ is satisfied by the initial
product-automaton synthesis and by each subsequent recomputation after
an event. Every confirmed contingent request is inserted into the
current task automaton as a dedicated layer and therefore becomes part
of the active specification. By
Theorem~\ref{thm:request_completion}, each such contingent task is
completed in finite time. Once completed, the corresponding layer is
removed, and the product automaton continues with the remaining active
tasks and the static suffix. Hence the static task, the contingent
service task, and the contingent formation task are all satisfied, which
implies the satisfaction of $\varphi_i$.
\end{proof}
